
\documentclass{article}
\usepackage{colortbl}
\usepackage{makecell}
\usepackage{multirow}
\usepackage{supertabular}

\begin{document}

\newcounter{utterance}

\centering \large Interaction Transcript for game `clean\_up', experiment `1\_medium\_5obj\_lv', episode 3 with mistral{-}large{-}3{-}azure.
\vspace{24pt}

{ \footnotesize  \setcounter{utterance}{1}
\setlength{\tabcolsep}{0pt}
\begin{supertabular}{c@{$\;$}|p{.15\linewidth}@{}p{.15\linewidth}p{.15\linewidth}p{.15\linewidth}p{.15\linewidth}p{.15\linewidth}}
    \# & $\;$A & \multicolumn{4}{c}{Game Master} & $\;\:$B\\
    \hline

    \theutterance \stepcounter{utterance}  
    & & \multicolumn{4}{p{0.6\linewidth}}{
        \cellcolor[rgb]{0.9,0.9,0.9}{
            \makecell[{{p{\linewidth}}}]{
                \texttt{\tiny{[P1$\langle$GM]}}
                \texttt{Es esmu jūsu spēles vadītājs, un jūs spēlējat sadarbības spēli ar šādu režģi kā spēles laukumu:} \\
\\ 
\texttt{```} \\
\texttt{ 1234567} \\
\texttt{╔═╤═════╗} \\
\texttt{║H│◌◌◌TI║ 1} \\
\texttt{╟─┤◌◌◌◌◌║ 2} \\
\texttt{║◌└─┐◌◌◌║ 3} \\
\texttt{║◌◌◌├─┐◌║ 4} \\
\texttt{║C◌◌│◌│◌║ 5} \\
\texttt{╟─┬─┴─┼─╢ 6} \\
\texttt{║◌│W◌◌│◌║ 7} \\
\texttt{╚═╧═══╧═╝} \\
\\ 
\texttt{```} \\
\\ 
\texttt{* Augšējā malā ir attēlotas x koordinātas, kas pieaug pa labi, un labajā malā y koordinātas, kas pieaug uz leju.} \\
\texttt{* Uz jūsu režģa ir nejauši izvietoti šādi objekti: 'W', 'I', 'T', 'C', 'H'.} \\
\\ 
\\ 
\texttt{Otrs spēlētājs redz spēles laukuma variāciju, kurā objekti ir izvietoti citās nejaušās vietās. Jūs nevarat redzēt otra spēlētāja laukumu, un viņš nevar redzēt jūsējo.} \\
\\ 
\\ 
\texttt{**Mērķis:**} \\
\\ 
\texttt{Abiem spēlētājiem ir jāpārvieto objekti uz sava fona tā, lai vienādi objekti nonāktu vienādās koordinātēs. Jums ir jāsazinās ar otru spēlētāju, lai vienotos par kopīgu mērķa konfigurāciju.} \\
\\ 
\\ 
\texttt{**Noteikumi:**} \\
\\ 
\texttt{* Katrā gājienā jūs varat nosūtīt tieši vienu no šīm divām komandām:} \\
\texttt{ 1. `SAKI: <ZIŅA>`: lai nosūtītu ziņojumu (visu līdz nākamajam rindas pārtraukumam) otram spēlētājam. Es to pārsūtīšu jūsu partnerim.} \\
\texttt{2. `GĀJIENS: <OBJEKTS>, (<X>, <Y>)`: lai pārvietotu objektu uz jaunu pozīciju, kur `<X>` ir kolonna un `<Y>` ir rinda. Es jūs informēšu, vai jūsu gājiens bija veiksmīgs vai nē.} \\
\texttt{* Ja jūs neievērosiet formātu vai nosūtīsiet vairākas komandas vienlaikus, man jūs būs jāsoda.} \\
\texttt{* Ja abi spēlētāji kopā uzkrāj vairāk nekā 12 sodus, jūs abi zaudējat spēli.} \\
\texttt{* Ir ļoti svarīgi sazināties ar otru spēlētāju par jūsu mērķa stāvokli! *Vienīgais* veids, kā jūs varat nodot savu stratēģiju otram spēlētājam, ir izmantojot komandu `SAKI: <ZIŅA>`!} \\
\\ 
\\ 
\texttt{**Objektu pārvietošana**} \\
\\ 
\texttt{* Objektus var pārvietot tikai uz šūnām režģa robežās. Mērķa šūnai jābūt tukšai, t.i., tajā jābūt tikai simbolam '◌'.} \\
\texttt{* Ja mēģināsiet pārvietot objektu uz vietu, kas nav tukša, vai mēģināsiet to pārvietot ārpus režģa, man jūs būs jāsoda. Jūs saņemat vēl vienu mēģinājumu.} \\
\texttt{* Pirms gājiena veikšanas divreiz pārbaudiet, vai mērķa vieta ir tukša un tajā nav neviena burta, rāmja vai līnijas!} \\
\\ 
\\ 
\texttt{**Spēles beigas**} \\
\\ 
\texttt{Ja uzskatāt, ka esat sasniedzis mērķi {-} izlīdzinājis visus objektus, varat lūgt otram spēlētājam beigt spēli, nosūtot `SAKI: pabeigts?`. Ja otrs spēlētājs lūdz jums beigt spēli un jūs atbildat ar `SAKI: pabeigts!`, spēle beigsies.} \\
\\ 
\texttt{Abi spēlētāji uzvar, ja spēle beidzas 20 raundos, kur viens raunds ir definēts kā divi spēlētāji, kas katrs nosūta derīgu komandu.} \\
\\ 
\texttt{**Punktu skaitīšana:**} \\
\\ 
\texttt{Jo tuvāk vienādi objekti atrodas abos spēles laukumos, jo vairāk punktu jūs saņemat. Sodi samazina jūsu punktus. Vai varat pārspēt rekordu?} \\
\\ 
\\ 
\texttt{Lūdzu, nosūtiet ziņu otram spēlētājam, lai sāktu spēli!} \\
            }
        }
    }
    & \\ \\

    \theutterance \stepcounter{utterance}  
    & \multicolumn{4}{p{0.6\linewidth}}{
        \cellcolor[rgb]{0.9,0.9,0.9}{
            \makecell[{{p{\linewidth}}}]{
                \texttt{\tiny{[P1$\rangle$GM]}}
                \texttt{Labi, sāksim spēli! Šeit ir mana pirmā ziņa, lai sāktu koordinēties ar tevi par mērķa konfigurāciju:} \\
\\ 
\texttt{```} \\
\texttt{SAKI: Sveiciens! Mani objekti ir šādās pozīcijās:} \\
\texttt{{-} H: (2, 1)} \\
\texttt{{-} T: (6, 1)} \\
\texttt{{-} I: (5, 1)} \\
\texttt{{-} C: (1, 5)} \\
\texttt{{-} W: (3, 7)} \\
\\ 
\texttt{Kādi ir tavi objekti? Mēs varam vienoties par mērķa koordinātām, piemēram, izvietojot objektus alfabēta secībā (C, H, I, T, W) vai kā citādi tev būtu ērti. Piedāvā savu variantu!} \\
\texttt{```} \\
\\ 
\texttt{Pēc šīs ziņas gaidīšu tavu atbildi, lai turpinātu plānot stratēģiju. Ja tev ir labāks priekšlikums par objektu izvietojumu, lūdzu, piedāvā to!} \\
            }
        }
    }
    & & \\ \\

    \theutterance \stepcounter{utterance}  
    & & & \multicolumn{2}{p{0.3\linewidth}}{
        \cellcolor[rgb]{0.9,0.9,0.9}{
            \makecell[{{p{\linewidth}}}]{
                \texttt{\tiny{[GM$|$GM]}}
                \texttt{Labi, sāksim spēli! Šeit ir mana pirmā ziņa, lai sāktu koordinēties ar tevi par mērķa konfigurāciju:} \\
\\ 
\texttt{```} \\
\texttt{SAKI: Sveiciens! Mani objekti ir šādās pozīcijās:} \\
\texttt{{-} H: (2, 1)} \\
\texttt{{-} T: (6, 1)} \\
\texttt{{-} I: (5, 1)} \\
\texttt{{-} C: (1, 5)} \\
\texttt{{-} W: (3, 7)} \\
\\ 
\texttt{Kādi ir tavi objekti? Mēs varam vienoties par mērķa koordinātām, piemēram, izvietojot objektus alfabēta secībā (C, H, I, T, W) vai kā citādi tev būtu ērti. Piedāvā savu variantu!} \\
\texttt{```} \\
\\ 
\texttt{Pēc šīs ziņas gaidīšu tavu atbildi, lai turpinātu plānot stratēģiju. Ja tev ir labāks priekšlikums par objektu izvietojumu, lūdzu, piedāvā to!} \\
            }
        }
    }
    & & \\ \\

    \theutterance \stepcounter{utterance}  
    & & & \multicolumn{2}{p{0.3\linewidth}}{
        \cellcolor[rgb]{0.9,0.9,0.9}{
            \makecell[{{p{\linewidth}}}]{
                \texttt{\tiny{[GM$|$GM]}}
                \texttt{Invalid format: head and tail are not empty} \\
\texttt{head: 'Labi, sāksim spēli! Šeit ir mana pirmā ziņa, lai sāktu koordinēties ar tevi par mērķa konfigurāciju:} \\
\\ 
\\ 
\texttt{'} \\
\texttt{tail: '} \\
\texttt{{-} H: (2, 1)} \\
\texttt{{-} T: (6, 1)} \\
\texttt{{-} I: (5, 1)} \\
\texttt{{-} C: (1, 5)} \\
\texttt{{-} W: (3, 7)} \\
\\ 
\texttt{Kādi ir tavi objekti? Mēs varam vienoties par mērķa koordinātām, piemēram, izvietojot objektus alfabēta secībā (C, H, I, T, W) vai kā citādi tev būtu ērti. Piedāvā savu variantu!} \\
\\ 
\\ 
\texttt{Pēc šīs ziņas gaidīšu tavu atbildi, lai turpinātu plānot stratēģiju. Ja tev ir labāks priekšlikums par objektu izvietojumu, lūdzu, piedāvā to!'} \\
            }
        }
    }
    & & \\ \\

    \theutterance \stepcounter{utterance}  
    & & \multicolumn{4}{p{0.6\linewidth}}{
        \cellcolor[rgb]{0.9,0.9,0.9}{
            \makecell[{{p{\linewidth}}}]{
                \texttt{\tiny{[P1$\langle$GM]}}
                \texttt{Sods: Jūsu ziņojumā nedrīkst būt nekas pirms vai pēc komandas!} \\
\\ 
\texttt{Pārliecinieties, ka jūsu atbildē ir tikai `SAKI: <ZIŅA>` vai `GĀJIENS: <OBJEKTS>, (<X>, <Y>)`, un nekas cits!} \\
\\ 
\texttt{Jūs kopā esat uzkrājuši 1 no 12 sodiem.} \\
\texttt{Lūdzu, mēģiniet vēlreiz!} \\
            }
        }
    }
    & \\ \\

    \theutterance \stepcounter{utterance}  
    & \multicolumn{4}{p{0.6\linewidth}}{
        \cellcolor[rgb]{0.9,0.9,0.9}{
            \makecell[{{p{\linewidth}}}]{
                \texttt{\tiny{[P1$\rangle$GM]}}
                \texttt{```} \\
\texttt{SAKI: Sveiciens! Mani objekti: H(2,1), T(6,1), I(5,1), C(1,5), W(3,7). Kādi ir tavi? Piedāvāsim mērķa koordinātas alfabēta secībā: C(1,1), H(2,1), I(3,1), T(4,1), W(5,1). Vai der?} \\
\texttt{```} \\
            }
        }
    }
    & & \\ \\

    \theutterance \stepcounter{utterance}  
    & & & \multicolumn{2}{p{0.3\linewidth}}{
        \cellcolor[rgb]{0.9,0.9,0.9}{
            \makecell[{{p{\linewidth}}}]{
                \texttt{\tiny{[GM$|$GM]}}
                \texttt{```} \\
\texttt{SAKI: Sveiciens! Mani objekti: H(2,1), T(6,1), I(5,1), C(1,5), W(3,7). Kādi ir tavi? Piedāvāsim mērķa koordinātas alfabēta secībā: C(1,1), H(2,1), I(3,1), T(4,1), W(5,1). Vai der?} \\
\texttt{```} \\
            }
        }
    }
    & & \\ \\

    \theutterance \stepcounter{utterance}  
    & & \multicolumn{4}{p{0.6\linewidth}}{
        \cellcolor[rgb]{0.9,0.9,0.9}{
            \makecell[{{p{\linewidth}}}]{
                \texttt{\tiny{[GM$\rangle$P2]}}
                \texttt{Es esmu jūsu spēles vadītājs, un jūs spēlējat sadarbības spēli ar šādu režģi kā spēles laukumu:} \\
\\ 
\texttt{```} \\
\texttt{ 1234567} \\
\texttt{╔═╤═════╗} \\
\texttt{║◌│◌◌◌◌◌║ 1} \\
\texttt{╟─┤◌◌◌H◌║ 2} \\
\texttt{║◌└─┐◌◌◌║ 3} \\
\texttt{║T◌I├─┐◌║ 4} \\
\texttt{║◌◌◌│C│◌║ 5} \\
\texttt{╟─┬─┴─┼─╢ 6} \\
\texttt{║◌│◌◌◌│W║ 7} \\
\texttt{╚═╧═══╧═╝} \\
\\ 
\texttt{```} \\
\\ 
\texttt{* Augšējā malā ir attēlotas x koordinātas, kas pieaug pa labi, un labajā malā y koordinātas, kas pieaug uz leju.} \\
\texttt{* Uz jūsu režģa ir nejauši izvietoti šādi objekti: 'W', 'I', 'T', 'C', 'H'.} \\
\\ 
\\ 
\texttt{Otrs spēlētājs redz spēles laukuma variāciju, kurā objekti ir izvietoti citās nejaušās vietās. Jūs nevarat redzēt otra spēlētāja laukumu, un viņš nevar redzēt jūsējo.} \\
\\ 
\\ 
\texttt{**Mērķis:**} \\
\\ 
\texttt{Abiem spēlētājiem ir jāpārvieto objekti uz sava fona tā, lai vienādi objekti nonāktu vienādās koordinātēs. Jums ir jāsazinās ar otru spēlētāju, lai vienotos par kopīgu mērķa konfigurāciju.} \\
\\ 
\\ 
\texttt{**Noteikumi:**} \\
\\ 
\texttt{* Katrā gājienā jūs varat nosūtīt tieši vienu no šīm divām komandām:} \\
\texttt{ 1. `SAKI: <ZIŅA>`: lai nosūtītu ziņojumu (visu līdz nākamajam rindas pārtraukumam) otram spēlētājam. Es to pārsūtīšu jūsu partnerim.} \\
\texttt{2. `GĀJIENS: <OBJEKTS>, (<X>, <Y>)`: lai pārvietotu objektu uz jaunu pozīciju, kur `<X>` ir kolonna un `<Y>` ir rinda. Es jūs informēšu, vai jūsu gājiens bija veiksmīgs vai nē.} \\
\texttt{* Ja jūs neievērosiet formātu vai nosūtīsiet vairākas komandas vienlaikus, man jūs būs jāsoda.} \\
\texttt{* Ja abi spēlētāji kopā uzkrāj vairāk nekā 12 sodus, jūs abi zaudējat spēli.} \\
\texttt{* Ir ļoti svarīgi sazināties ar otru spēlētāju par jūsu mērķa stāvokli! *Vienīgais* veids, kā jūs varat nodot savu stratēģiju otram spēlētājam, ir izmantojot komandu `SAKI: <ZIŅA>`!} \\
\\ 
\\ 
\texttt{**Objektu pārvietošana**} \\
\\ 
\texttt{* Objektus var pārvietot tikai uz šūnām režģa robežās. Mērķa šūnai jābūt tukšai, t.i., tajā jābūt tikai simbolam '◌'.} \\
\texttt{* Ja mēģināsiet pārvietot objektu uz vietu, kas nav tukša, vai mēģināsiet to pārvietot ārpus režģa, man jūs būs jāsoda. Jūs saņemat vēl vienu mēģinājumu.} \\
\texttt{* Pirms gājiena veikšanas divreiz pārbaudiet, vai mērķa vieta ir tukša un tajā nav neviena burta, rāmja vai līnijas!} \\
\\ 
\\ 
\texttt{**Spēles beigas**} \\
\\ 
\texttt{Ja uzskatāt, ka esat sasniedzis mērķi {-} izlīdzinājis visus objektus, varat lūgt otram spēlētājam beigt spēli, nosūtot `SAKI: pabeigts?`. Ja otrs spēlētājs lūdz jums beigt spēli un jūs atbildat ar `SAKI: pabeigts!`, spēle beigsies.} \\
\\ 
\texttt{Abi spēlētāji uzvar, ja spēle beidzas 20 raundos, kur viens raunds ir definēts kā divi spēlētāji, kas katrs nosūta derīgu komandu.} \\
\\ 
\texttt{**Punktu skaitīšana:**} \\
\\ 
\texttt{Jo tuvāk vienādi objekti atrodas abos spēles laukumos, jo vairāk punktu jūs saņemat. Sodi samazina jūsu punktus. Vai varat pārspēt rekordu?} \\
\\ 
\texttt{Otrs spēlētājs sāka spēli, nosūtot šo ziņojumu:} \\
\\ 
\texttt{"Sveiciens! Mani objekti: H(2,1), T(6,1), I(5,1), C(1,5), W(3,7). Kādi ir tavi? Piedāvāsim mērķa koordinātas alfabēta secībā: C(1,1), H(2,1), I(3,1), T(4,1), W(5,1). Vai der?"} \\
\\ 
\texttt{Kāda ir jūsu pirmā komanda?} \\
            }
        }
    }
    & \\ \\

    \theutterance \stepcounter{utterance}  
    & & & \multicolumn{4}{p{0.6\linewidth}}{
        \cellcolor[rgb]{0.9,0.9,0.9}{
            \makecell[{{p{\linewidth}}}]{
                \texttt{\tiny{[GM$\langle$P2]}}
                \texttt{Pēc otra spēlētāja ziņas analīzes, man ir skaidrs, ka viņš piedāvā mērķa koordinātas alfabēta secībā (C, H, I, T, W) augšējā rindā (y=1). Lai nodrošinātu sadarbību, man ir jāpaziņo savi objekti un jāpiekrīt vai jāpiedāvā alternatīva.} \\
\\ 
\texttt{Mana pašreizējā konfigurācija ir:} \\
\texttt{{-} H(5,2)} \\
\texttt{{-} T(1,4)} \\
\texttt{{-} I(3,4)} \\
\texttt{{-} C(4,5)} \\
\texttt{{-} W(7,7)} \\
\\ 
\texttt{\#\#\# Pirmā komanda:} \\
\texttt{Es sākšu ar ziņojumu, lai apstiprinātu savus objektus un piedāvāto mērķi, kā arī piedāvāšu nelielu korekciju, lai izvairītos no konfliktiem ar esošajiem objektiem (piemēram, H jau atrodas (2,1) otram spēlētājam, bet man tas ir (5,2), tāpēc alfabēta secība nav pilnīgi iespējama bez pārvietošanas).} \\
\\ 
\texttt{```} \\
\texttt{SAKI: Sveiciens! Mani objekti: H(5,2), T(1,4), I(3,4), C(4,5), W(7,7). Tava piedāvātā mērķa konfigurācija ir loģiska, bet man H nav (2,1), tāpēc piedāvāju šādu mērķi: C(1,1), H(2,2), I(3,1), T(4,1), W(5,1). Vai der?} \\
\texttt{```} \\
\\ 
\texttt{Šis piedāvājums:} \\
\texttt{1. Saglabā alfabēta secību (C, H, I, T, W).} \\
\texttt{2. Izvietojumā izvairās no pārklāšanās ar esošajiem objektiem (piemēram, H(2,2) ir tuvu manam H(5,2), bet pietiekami tālu, lai izvairītos no sākotnējām kļūdām).} \\
\texttt{3. Ļauj abiem spēlētājiem pakāpeniski pārvietot objektus uz mērķa pozīcijām.} \\
\\ 
\texttt{Nākamais solis būs gaidīt otra spēlētāja atbildi un sākt pārvietošanas gājienus.} \\
            }
        }
    }
     \\ \\

    \theutterance \stepcounter{utterance}  
    & & & \multicolumn{2}{p{0.3\linewidth}}{
        \cellcolor[rgb]{0.9,0.9,0.9}{
            \makecell[{{p{\linewidth}}}]{
                \texttt{\tiny{[GM$|$GM]}}
                \texttt{Pēc otra spēlētāja ziņas analīzes, man ir skaidrs, ka viņš piedāvā mērķa koordinātas alfabēta secībā (C, H, I, T, W) augšējā rindā (y=1). Lai nodrošinātu sadarbību, man ir jāpaziņo savi objekti un jāpiekrīt vai jāpiedāvā alternatīva.} \\
\\ 
\texttt{Mana pašreizējā konfigurācija ir:} \\
\texttt{{-} H(5,2)} \\
\texttt{{-} T(1,4)} \\
\texttt{{-} I(3,4)} \\
\texttt{{-} C(4,5)} \\
\texttt{{-} W(7,7)} \\
\\ 
\texttt{\#\#\# Pirmā komanda:} \\
\texttt{Es sākšu ar ziņojumu, lai apstiprinātu savus objektus un piedāvāto mērķi, kā arī piedāvāšu nelielu korekciju, lai izvairītos no konfliktiem ar esošajiem objektiem (piemēram, H jau atrodas (2,1) otram spēlētājam, bet man tas ir (5,2), tāpēc alfabēta secība nav pilnīgi iespējama bez pārvietošanas).} \\
\\ 
\texttt{```} \\
\texttt{SAKI: Sveiciens! Mani objekti: H(5,2), T(1,4), I(3,4), C(4,5), W(7,7). Tava piedāvātā mērķa konfigurācija ir loģiska, bet man H nav (2,1), tāpēc piedāvāju šādu mērķi: C(1,1), H(2,2), I(3,1), T(4,1), W(5,1). Vai der?} \\
\texttt{```} \\
\\ 
\texttt{Šis piedāvājums:} \\
\texttt{1. Saglabā alfabēta secību (C, H, I, T, W).} \\
\texttt{2. Izvietojumā izvairās no pārklāšanās ar esošajiem objektiem (piemēram, H(2,2) ir tuvu manam H(5,2), bet pietiekami tālu, lai izvairītos no sākotnējām kļūdām).} \\
\texttt{3. Ļauj abiem spēlētājiem pakāpeniski pārvietot objektus uz mērķa pozīcijām.} \\
\\ 
\texttt{Nākamais solis būs gaidīt otra spēlētāja atbildi un sākt pārvietošanas gājienus.} \\
            }
        }
    }
    & & \\ \\

    \theutterance \stepcounter{utterance}  
    & & & \multicolumn{2}{p{0.3\linewidth}}{
        \cellcolor[rgb]{0.9,0.9,0.9}{
            \makecell[{{p{\linewidth}}}]{
                \texttt{\tiny{[GM$|$GM]}}
                \texttt{Invalid format: head and tail are not empty} \\
\texttt{head: 'Pēc otra spēlētāja ziņas analīzes, man ir skaidrs, ka viņš piedāvā mērķa koordinātas alfabēta secībā (C, H, I, T, W) augšējā rindā (y=1). Lai nodrošinātu sadarbību, man ir jāpaziņo savi objekti un jāpiekrīt vai jāpiedāvā alternatīva.} \\
\\ 
\texttt{Mana pašreizējā konfigurācija ir:} \\
\texttt{{-} H(5,2)} \\
\texttt{{-} T(1,4)} \\
\texttt{{-} I(3,4)} \\
\texttt{{-} C(4,5)} \\
\texttt{{-} W(7,7)} \\
\\ 
\texttt{\#\#\# Pirmā komanda:} \\
\texttt{Es sākšu ar ziņojumu, lai apstiprinātu savus objektus un piedāvāto mērķi, kā arī piedāvāšu nelielu korekciju, lai izvairītos no konfliktiem ar esošajiem objektiem (piemēram, H jau atrodas (2,1) otram spēlētājam, bet man tas ir (5,2), tāpēc alfabēta secība nav pilnīgi iespējama bez pārvietošanas).} \\
\\ 
\\ 
\texttt{'} \\
\texttt{tail: '} \\
\\ 
\\ 
\texttt{Šis piedāvājums:} \\
\texttt{1. Saglabā alfabēta secību (C, H, I, T, W).} \\
\texttt{2. Izvietojumā izvairās no pārklāšanās ar esošajiem objektiem (piemēram, H(2,2) ir tuvu manam H(5,2), bet pietiekami tālu, lai izvairītos no sākotnējām kļūdām).} \\
\texttt{3. Ļauj abiem spēlētājiem pakāpeniski pārvietot objektus uz mērķa pozīcijām.} \\
\\ 
\texttt{Nākamais solis būs gaidīt otra spēlētāja atbildi un sākt pārvietošanas gājienus.'} \\
            }
        }
    }
    & & \\ \\

    \theutterance \stepcounter{utterance}  
    & & \multicolumn{4}{p{0.6\linewidth}}{
        \cellcolor[rgb]{0.9,0.9,0.9}{
            \makecell[{{p{\linewidth}}}]{
                \texttt{\tiny{[GM$\rangle$P2]}}
                \texttt{Sods: Jūsu ziņojumā nedrīkst būt nekas pirms vai pēc komandas!} \\
\\ 
\texttt{Pārliecinieties, ka jūsu atbildē ir tikai `SAKI: <ZIŅA>` vai `GĀJIENS: <OBJEKTS>, (<X>, <Y>)`, un nekas cits!} \\
\\ 
\texttt{Jūs kopā esat uzkrājuši 2 no 12 sodiem.} \\
\texttt{Lūdzu, mēģiniet vēlreiz!} \\
            }
        }
    }
    & \\ \\

    \theutterance \stepcounter{utterance}  
    & & & \multicolumn{4}{p{0.6\linewidth}}{
        \cellcolor[rgb]{0.9,0.9,0.9}{
            \makecell[{{p{\linewidth}}}]{
                \texttt{\tiny{[GM$\langle$P2]}}
                \texttt{Atvainojos par kļūdu! Šoreiz ievērošu formātu precīzi.} \\
\\ 
\texttt{Pirmā komanda (tikai ziņojums bez liekas informācijas):} \\
\texttt{```} \\
\texttt{SAKI: Mani objekti: H(5,2), T(1,4), I(3,4), C(4,5), W(7,7). Piedāvāju mērķi: C(1,1), H(2,2), I(3,1), T(4,1), W(5,1). Vai der?} \\
\texttt{```} \\
            }
        }
    }
     \\ \\

    \theutterance \stepcounter{utterance}  
    & & & \multicolumn{2}{p{0.3\linewidth}}{
        \cellcolor[rgb]{0.9,0.9,0.9}{
            \makecell[{{p{\linewidth}}}]{
                \texttt{\tiny{[GM$|$GM]}}
                \texttt{Atvainojos par kļūdu! Šoreiz ievērošu formātu precīzi.} \\
\\ 
\texttt{Pirmā komanda (tikai ziņojums bez liekas informācijas):} \\
\texttt{```} \\
\texttt{SAKI: Mani objekti: H(5,2), T(1,4), I(3,4), C(4,5), W(7,7). Piedāvāju mērķi: C(1,1), H(2,2), I(3,1), T(4,1), W(5,1). Vai der?} \\
\texttt{```} \\
            }
        }
    }
    & & \\ \\

    \theutterance \stepcounter{utterance}  
    & & & \multicolumn{2}{p{0.3\linewidth}}{
        \cellcolor[rgb]{0.9,0.9,0.9}{
            \makecell[{{p{\linewidth}}}]{
                \texttt{\tiny{[GM$|$GM]}}
                \texttt{Invalid format: head is not empty: 'Atvainojos par kļūdu! Šoreiz ievērošu formātu precīzi.} \\
\\ 
\texttt{Pirmā komanda (tikai ziņojums bez liekas informācijas):} \\
\\ 
\texttt{'} \\
            }
        }
    }
    & & \\ \\

    \theutterance \stepcounter{utterance}  
    & & \multicolumn{4}{p{0.6\linewidth}}{
        \cellcolor[rgb]{0.9,0.9,0.9}{
            \makecell[{{p{\linewidth}}}]{
                \texttt{\tiny{[GM$\rangle$P2]}}
                \texttt{Sods: Jūsu ziņojumā nedrīkst būt nekas pirms komandas!} \\
\\ 
\texttt{Pārliecinieties, ka jūsu atbildē ir tikai `SAKI: <ZIŅA>` vai `GĀJIENS: <OBJEKTS>, (<X>, <Y>)`, un nekas cits!} \\
\\ 
\texttt{Jūs kopā esat uzkrājuši 3 no 12 sodiem.} \\
\texttt{Lūdzu, mēģiniet vēlreiz!} \\
            }
        }
    }
    & \\ \\

    \theutterance \stepcounter{utterance}  
    & & & \multicolumn{4}{p{0.6\linewidth}}{
        \cellcolor[rgb]{0.9,0.9,0.9}{
            \makecell[{{p{\linewidth}}}]{
                \texttt{\tiny{[GM$\langle$P2]}}
                \texttt{Atvainojos par iepriekšējām kļūdām. Tagad sūtīšu tikai precīzu ziņojumu bez jebkādām papildu frāzēm:} \\
\\ 
\texttt{```} \\
\texttt{SAKI: H(5,2) T(1,4) I(3,4) C(4,5) W(7,7). Mērķis: C(1,1) H(2,2) I(3,1) T(4,1) W(5,1). Der?} \\
\texttt{```} \\
            }
        }
    }
     \\ \\

    \theutterance \stepcounter{utterance}  
    & & & \multicolumn{2}{p{0.3\linewidth}}{
        \cellcolor[rgb]{0.9,0.9,0.9}{
            \makecell[{{p{\linewidth}}}]{
                \texttt{\tiny{[GM$|$GM]}}
                \texttt{Atvainojos par iepriekšējām kļūdām. Tagad sūtīšu tikai precīzu ziņojumu bez jebkādām papildu frāzēm:} \\
\\ 
\texttt{```} \\
\texttt{SAKI: H(5,2) T(1,4) I(3,4) C(4,5) W(7,7). Mērķis: C(1,1) H(2,2) I(3,1) T(4,1) W(5,1). Der?} \\
\texttt{```} \\
            }
        }
    }
    & & \\ \\

    \theutterance \stepcounter{utterance}  
    & & & \multicolumn{2}{p{0.3\linewidth}}{
        \cellcolor[rgb]{0.9,0.9,0.9}{
            \makecell[{{p{\linewidth}}}]{
                \texttt{\tiny{[GM$|$GM]}}
                \texttt{Invalid format: head is not empty: 'Atvainojos par iepriekšējām kļūdām. Tagad sūtīšu tikai precīzu ziņojumu bez jebkādām papildu frāzēm:} \\
\\ 
\\ 
\texttt{'} \\
            }
        }
    }
    & & \\ \\

    \theutterance \stepcounter{utterance}  
    & & \multicolumn{4}{p{0.6\linewidth}}{
        \cellcolor[rgb]{0.9,0.9,0.9}{
            \makecell[{{p{\linewidth}}}]{
                \texttt{\tiny{[GM$\rangle$P2]}}
                \texttt{Sods: Jūsu ziņojumā nedrīkst būt nekas pirms komandas!} \\
\\ 
\texttt{Pārliecinieties, ka jūsu atbildē ir tikai `SAKI: <ZIŅA>` vai `GĀJIENS: <OBJEKTS>, (<X>, <Y>)`, un nekas cits!} \\
\\ 
\texttt{Jūs kopā esat uzkrājuši 4 no 12 sodiem.} \\
\texttt{Lūdzu, mēģiniet vēlreiz!} \\
            }
        }
    }
    & \\ \\

    \theutterance \stepcounter{utterance}  
    & & & \multicolumn{4}{p{0.6\linewidth}}{
        \cellcolor[rgb]{0.9,0.9,0.9}{
            \makecell[{{p{\linewidth}}}]{
                \texttt{\tiny{[GM$\langle$P2]}}
                \texttt{Tagad esmu sapratis {-} ziņojumam jābūt **tikai** formātā `SAKI: <teksts>` bez papildu komentāriem. Mēģināšu vēlreiz:} \\
\\ 
\texttt{```} \\
\texttt{SAKI: Mani objekti: H(5,2), T(1,4), I(3,4), C(4,5), W(7,7). Piedāvāju mērķus alfabētiski: C(1,1), H(2,2), I(3,1), T(4,1), W(5,1).} \\
\texttt{```} \\
            }
        }
    }
     \\ \\

    \theutterance \stepcounter{utterance}  
    & & & \multicolumn{2}{p{0.3\linewidth}}{
        \cellcolor[rgb]{0.9,0.9,0.9}{
            \makecell[{{p{\linewidth}}}]{
                \texttt{\tiny{[GM$|$GM]}}
                \texttt{Tagad esmu sapratis {-} ziņojumam jābūt **tikai** formātā `SAKI: <teksts>` bez papildu komentāriem. Mēģināšu vēlreiz:} \\
\\ 
\texttt{```} \\
\texttt{SAKI: Mani objekti: H(5,2), T(1,4), I(3,4), C(4,5), W(7,7). Piedāvāju mērķus alfabētiski: C(1,1), H(2,2), I(3,1), T(4,1), W(5,1).} \\
\texttt{```} \\
            }
        }
    }
    & & \\ \\

    \theutterance \stepcounter{utterance}  
    & & & \multicolumn{2}{p{0.3\linewidth}}{
        \cellcolor[rgb]{0.9,0.9,0.9}{
            \makecell[{{p{\linewidth}}}]{
                \texttt{\tiny{[GM$|$GM]}}
                \texttt{Invalid format: head is not empty: 'Tagad esmu sapratis {-} ziņojumam jābūt **tikai** formātā SAKI: <teksts> bez papildu komentāriem. Mēģināšu vēlreiz:} \\
\\ 
\\ 
\texttt{'} \\
            }
        }
    }
    & & \\ \\

    \theutterance \stepcounter{utterance}  
    & & \multicolumn{4}{p{0.6\linewidth}}{
        \cellcolor[rgb]{0.9,0.9,0.9}{
            \makecell[{{p{\linewidth}}}]{
                \texttt{\tiny{[GM$\rangle$P2]}}
                \texttt{Sods: Jūsu ziņojumā nedrīkst būt nekas pirms komandas!} \\
\\ 
\texttt{Pārliecinieties, ka jūsu atbildē ir tikai `SAKI: <ZIŅA>` vai `GĀJIENS: <OBJEKTS>, (<X>, <Y>)`, un nekas cits!} \\
\\ 
\texttt{Jūs kopā esat uzkrājuši 5 no 12 sodiem.} \\
\texttt{Lūdzu, mēģiniet vēlreiz!} \\
            }
        }
    }
    & \\ \\

    \theutterance \stepcounter{utterance}  
    & & & \multicolumn{4}{p{0.6\linewidth}}{
        \cellcolor[rgb]{0.9,0.9,0.9}{
            \makecell[{{p{\linewidth}}}]{
                \texttt{\tiny{[GM$\langle$P2]}}
                \texttt{Tagad saprotu {-} man jāizdrukā **tikai** komandas teksts bez jebkādiem ievada aprakstiem. Šeit ir pareizs formāts:} \\
\\ 
\texttt{```} \\
\texttt{SAKI: H(5,2) T(1,4) I(3,4) C(4,5) W(7,7). Mērķi: C(1,1) H(2,2) I(3,1) T(4,1) W(5,1)} \\
\texttt{```} \\
            }
        }
    }
     \\ \\

    \theutterance \stepcounter{utterance}  
    & & & \multicolumn{2}{p{0.3\linewidth}}{
        \cellcolor[rgb]{0.9,0.9,0.9}{
            \makecell[{{p{\linewidth}}}]{
                \texttt{\tiny{[GM$|$GM]}}
                \texttt{Tagad saprotu {-} man jāizdrukā **tikai** komandas teksts bez jebkādiem ievada aprakstiem. Šeit ir pareizs formāts:} \\
\\ 
\texttt{```} \\
\texttt{SAKI: H(5,2) T(1,4) I(3,4) C(4,5) W(7,7). Mērķi: C(1,1) H(2,2) I(3,1) T(4,1) W(5,1)} \\
\texttt{```} \\
            }
        }
    }
    & & \\ \\

    \theutterance \stepcounter{utterance}  
    & & & \multicolumn{2}{p{0.3\linewidth}}{
        \cellcolor[rgb]{0.9,0.9,0.9}{
            \makecell[{{p{\linewidth}}}]{
                \texttt{\tiny{[GM$|$GM]}}
                \texttt{Invalid format: head is not empty: 'Tagad saprotu {-} man jāizdrukā **tikai** komandas teksts bez jebkādiem ievada aprakstiem. Šeit ir pareizs formāts:} \\
\\ 
\\ 
\texttt{'} \\
            }
        }
    }
    & & \\ \\

    \theutterance \stepcounter{utterance}  
    & & \multicolumn{4}{p{0.6\linewidth}}{
        \cellcolor[rgb]{0.9,0.9,0.9}{
            \makecell[{{p{\linewidth}}}]{
                \texttt{\tiny{[GM$\rangle$P2]}}
                \texttt{Sods: Jūsu ziņojumā nedrīkst būt nekas pirms komandas!} \\
\\ 
\texttt{Pārliecinieties, ka jūsu atbildē ir tikai `SAKI: <ZIŅA>` vai `GĀJIENS: <OBJEKTS>, (<X>, <Y>)`, un nekas cits!} \\
\\ 
\texttt{Jūs kopā esat uzkrājuši 6 no 12 sodiem.} \\
\texttt{Lūdzu, mēģiniet vēlreiz!} \\
            }
        }
    }
    & \\ \\

    \theutterance \stepcounter{utterance}  
    & & & \multicolumn{4}{p{0.6\linewidth}}{
        \cellcolor[rgb]{0.9,0.9,0.9}{
            \makecell[{{p{\linewidth}}}]{
                \texttt{\tiny{[GM$\langle$P2]}}
                \texttt{Atvainojos par nepārtrauktām kļūdām. Tagad sūtīšu tikai precīzu komandu bez jebkādiem papildu vārdiem:} \\
\\ 
\texttt{```} \\
\texttt{SAKI: H(5,2), T(1,4), I(3,4), C(4,5), W(7,7). Mērķi: C(1,1), H(2,2), I(3,1), T(4,1), W(5,1)} \\
\texttt{```} \\
            }
        }
    }
     \\ \\

    \theutterance \stepcounter{utterance}  
    & & & \multicolumn{2}{p{0.3\linewidth}}{
        \cellcolor[rgb]{0.9,0.9,0.9}{
            \makecell[{{p{\linewidth}}}]{
                \texttt{\tiny{[GM$|$GM]}}
                \texttt{Atvainojos par nepārtrauktām kļūdām. Tagad sūtīšu tikai precīzu komandu bez jebkādiem papildu vārdiem:} \\
\\ 
\texttt{```} \\
\texttt{SAKI: H(5,2), T(1,4), I(3,4), C(4,5), W(7,7). Mērķi: C(1,1), H(2,2), I(3,1), T(4,1), W(5,1)} \\
\texttt{```} \\
            }
        }
    }
    & & \\ \\

    \theutterance \stepcounter{utterance}  
    & & & \multicolumn{2}{p{0.3\linewidth}}{
        \cellcolor[rgb]{0.9,0.9,0.9}{
            \makecell[{{p{\linewidth}}}]{
                \texttt{\tiny{[GM$|$GM]}}
                \texttt{Invalid format: head is not empty: 'Atvainojos par nepārtrauktām kļūdām. Tagad sūtīšu tikai precīzu komandu bez jebkādiem papildu vārdiem:} \\
\\ 
\\ 
\texttt{'} \\
            }
        }
    }
    & & \\ \\

    \theutterance \stepcounter{utterance}  
    & & \multicolumn{4}{p{0.6\linewidth}}{
        \cellcolor[rgb]{0.9,0.9,0.9}{
            \makecell[{{p{\linewidth}}}]{
                \texttt{\tiny{[GM$\rangle$P2]}}
                \texttt{Sods: Jūsu ziņojumā nedrīkst būt nekas pirms komandas!} \\
\\ 
\texttt{Pārliecinieties, ka jūsu atbildē ir tikai `SAKI: <ZIŅA>` vai `GĀJIENS: <OBJEKTS>, (<X>, <Y>)`, un nekas cits!} \\
\\ 
\texttt{Jūs kopā esat uzkrājuši 7 no 12 sodiem.} \\
\texttt{Lūdzu, mēģiniet vēlreiz!} \\
            }
        }
    }
    & \\ \\

    \theutterance \stepcounter{utterance}  
    & & & \multicolumn{4}{p{0.6\linewidth}}{
        \cellcolor[rgb]{0.9,0.9,0.9}{
            \makecell[{{p{\linewidth}}}]{
                \texttt{\tiny{[GM$\langle$P2]}}
                \texttt{Es pilnībā saprotu noteikumus tagad. Šeit ir precīzi formāta ziņojums bez liekiem vārdiem:} \\
\\ 
\texttt{```} \\
\texttt{SAKI: Savs: H52 T14 I34 C45 W77. Mērķis: C11 H22 I31 T41 W51} \\
\texttt{```} \\
            }
        }
    }
     \\ \\

    \theutterance \stepcounter{utterance}  
    & & & \multicolumn{2}{p{0.3\linewidth}}{
        \cellcolor[rgb]{0.9,0.9,0.9}{
            \makecell[{{p{\linewidth}}}]{
                \texttt{\tiny{[GM$|$GM]}}
                \texttt{Es pilnībā saprotu noteikumus tagad. Šeit ir precīzi formāta ziņojums bez liekiem vārdiem:} \\
\\ 
\texttt{```} \\
\texttt{SAKI: Savs: H52 T14 I34 C45 W77. Mērķis: C11 H22 I31 T41 W51} \\
\texttt{```} \\
            }
        }
    }
    & & \\ \\

    \theutterance \stepcounter{utterance}  
    & & & \multicolumn{2}{p{0.3\linewidth}}{
        \cellcolor[rgb]{0.9,0.9,0.9}{
            \makecell[{{p{\linewidth}}}]{
                \texttt{\tiny{[GM$|$GM]}}
                \texttt{Invalid format: head is not empty: 'Es pilnībā saprotu noteikumus tagad. Šeit ir precīzi formāta ziņojums bez liekiem vārdiem:} \\
\\ 
\\ 
\texttt{'} \\
            }
        }
    }
    & & \\ \\

    \theutterance \stepcounter{utterance}  
    & & \multicolumn{4}{p{0.6\linewidth}}{
        \cellcolor[rgb]{0.9,0.9,0.9}{
            \makecell[{{p{\linewidth}}}]{
                \texttt{\tiny{[GM$\rangle$P2]}}
                \texttt{Sods: Jūsu ziņojumā nedrīkst būt nekas pirms komandas!} \\
\\ 
\texttt{Pārliecinieties, ka jūsu atbildē ir tikai `SAKI: <ZIŅA>` vai `GĀJIENS: <OBJEKTS>, (<X>, <Y>)`, un nekas cits!} \\
\\ 
\texttt{Jūs kopā esat uzkrājuši 8 no 12 sodiem.} \\
\texttt{Lūdzu, mēģiniet vēlreiz!} \\
            }
        }
    }
    & \\ \\

    \theutterance \stepcounter{utterance}  
    & & & \multicolumn{4}{p{0.6\linewidth}}{
        \cellcolor[rgb]{0.9,0.9,0.9}{
            \makecell[{{p{\linewidth}}}]{
                \texttt{\tiny{[GM$\langle$P2]}}
                \texttt{Tagad esmu pilnībā sapratis {-} ziņojumam jābūt tikai šādā formātā bez jebkādām papildu frāzēm:} \\
\\ 
\texttt{```} \\
\texttt{SAKI: H(5,2) T(1,4) I(3,4) C(4,5) W(7,7). C(1,1) H(2,2) I(3,1) T(4,1) W(5,1)} \\
\texttt{```} \\
\\ 
\texttt{(Nepievienoju vārdus "Savs" vai "Mērķis", tikai koordinātes)} \\
            }
        }
    }
     \\ \\

    \theutterance \stepcounter{utterance}  
    & & & \multicolumn{2}{p{0.3\linewidth}}{
        \cellcolor[rgb]{0.9,0.9,0.9}{
            \makecell[{{p{\linewidth}}}]{
                \texttt{\tiny{[GM$|$GM]}}
                \texttt{Tagad esmu pilnībā sapratis {-} ziņojumam jābūt tikai šādā formātā bez jebkādām papildu frāzēm:} \\
\\ 
\texttt{```} \\
\texttt{SAKI: H(5,2) T(1,4) I(3,4) C(4,5) W(7,7). C(1,1) H(2,2) I(3,1) T(4,1) W(5,1)} \\
\texttt{```} \\
\\ 
\texttt{(Nepievienoju vārdus "Savs" vai "Mērķis", tikai koordinātes)} \\
            }
        }
    }
    & & \\ \\

    \theutterance \stepcounter{utterance}  
    & & & \multicolumn{2}{p{0.3\linewidth}}{
        \cellcolor[rgb]{0.9,0.9,0.9}{
            \makecell[{{p{\linewidth}}}]{
                \texttt{\tiny{[GM$|$GM]}}
                \texttt{Invalid format: head and tail are not empty} \\
\texttt{head: 'Tagad esmu pilnībā sapratis {-} ziņojumam jābūt tikai šādā formātā bez jebkādām papildu frāzēm:} \\
\\ 
\\ 
\texttt{'} \\
\texttt{tail: '} \\
\\ 
\\ 
\texttt{(Nepievienoju vārdus "Savs" vai "Mērķis", tikai koordinātes)'} \\
            }
        }
    }
    & & \\ \\

    \theutterance \stepcounter{utterance}  
    & & \multicolumn{4}{p{0.6\linewidth}}{
        \cellcolor[rgb]{0.9,0.9,0.9}{
            \makecell[{{p{\linewidth}}}]{
                \texttt{\tiny{[GM$\rangle$P2]}}
                \texttt{Sods: Jūsu ziņojumā nedrīkst būt nekas pirms vai pēc komandas!} \\
\\ 
\texttt{Pārliecinieties, ka jūsu atbildē ir tikai `SAKI: <ZIŅA>` vai `GĀJIENS: <OBJEKTS>, (<X>, <Y>)`, un nekas cits!} \\
\\ 
\texttt{Jūs kopā esat uzkrājuši 9 no 12 sodiem.} \\
\texttt{Lūdzu, mēģiniet vēlreiz!} \\
            }
        }
    }
    & \\ \\

    \theutterance \stepcounter{utterance}  
    & & & \multicolumn{4}{p{0.6\linewidth}}{
        \cellcolor[rgb]{0.9,0.9,0.9}{
            \makecell[{{p{\linewidth}}}]{
                \texttt{\tiny{[GM$\langle$P2]}}
                \texttt{Es saprotu problēmu {-} man jāizdrukā tikai precīzi formāta ziņojums bez jebkādām papildu zīmēm vai atstarpēm. Šeit ir pareizais formāts:} \\
\\ 
\texttt{```} \\
\texttt{SAKI:H(5,2)T(1,4)I(3,4)C(4,5)W(7,7)Cilj:C(1,1)H(2,2)I(3,1)T(4,1)W(5,1)} \\
\texttt{```} \\
            }
        }
    }
     \\ \\

    \theutterance \stepcounter{utterance}  
    & & & \multicolumn{2}{p{0.3\linewidth}}{
        \cellcolor[rgb]{0.9,0.9,0.9}{
            \makecell[{{p{\linewidth}}}]{
                \texttt{\tiny{[GM$|$GM]}}
                \texttt{Es saprotu problēmu {-} man jāizdrukā tikai precīzi formāta ziņojums bez jebkādām papildu zīmēm vai atstarpēm. Šeit ir pareizais formāts:} \\
\\ 
\texttt{```} \\
\texttt{SAKI:H(5,2)T(1,4)I(3,4)C(4,5)W(7,7)Cilj:C(1,1)H(2,2)I(3,1)T(4,1)W(5,1)} \\
\texttt{```} \\
            }
        }
    }
    & & \\ \\

    \theutterance \stepcounter{utterance}  
    & & & \multicolumn{2}{p{0.3\linewidth}}{
        \cellcolor[rgb]{0.9,0.9,0.9}{
            \makecell[{{p{\linewidth}}}]{
                \texttt{\tiny{[GM$|$GM]}}
                \texttt{Invalid response format} \\
            }
        }
    }
    & & \\ \\

    \theutterance \stepcounter{utterance}  
    & & \multicolumn{4}{p{0.6\linewidth}}{
        \cellcolor[rgb]{0.9,0.9,0.9}{
            \makecell[{{p{\linewidth}}}]{
                \texttt{\tiny{[GM$\rangle$P2]}}
                \texttt{Sods: Jūsu ziņojums nav gaidītajā formātā!} \\
\\ 
\texttt{Pārliecinieties, ka jūsu atbildē ir tikai `SAKI: <ZIŅA>` vai `GĀJIENS: <OBJEKTS>, (<X>, <Y>)`, un nekas cits!} \\
\\ 
\texttt{Jūs kopā esat uzkrājuši 10 no 12 sodiem.} \\
\texttt{Lūdzu, mēģiniet vēlreiz!} \\
            }
        }
    }
    & \\ \\

    \theutterance \stepcounter{utterance}  
    & & & \multicolumn{4}{p{0.6\linewidth}}{
        \cellcolor[rgb]{0.9,0.9,0.9}{
            \makecell[{{p{\linewidth}}}]{
                \texttt{\tiny{[GM$\langle$P2]}}
                \texttt{Atvainojos par iepriekšējām kļūdām. Tagad es sūtīšu tikai divas atsevišķas ziņas precīzā formātā:} \\
\\ 
\texttt{Pirmā ziņa:} \\
\texttt{```} \\
\texttt{SAKI: Mani objekti: H(5,2) T(1,4) I(3,4) C(4,5) W(7,7)} \\
\texttt{```} \\
\\ 
\texttt{Otrā ziņa (ja pirmā tiks pieņemta):} \\
\texttt{```} \\
\texttt{SAKI: Piedāvātais mērķis: C(1,1) H(2,2) I(3,1) T(4,1) W(5,1)} \\
\texttt{```} \\
\\ 
\texttt{Vai vēlaties, lai es sūtu tikai pirmo ziņu? (Mēs esam tuvu maksimālajam sodu skaitam)} \\
            }
        }
    }
     \\ \\

    \theutterance \stepcounter{utterance}  
    & & & \multicolumn{2}{p{0.3\linewidth}}{
        \cellcolor[rgb]{0.9,0.9,0.9}{
            \makecell[{{p{\linewidth}}}]{
                \texttt{\tiny{[GM$|$GM]}}
                \texttt{Atvainojos par iepriekšējām kļūdām. Tagad es sūtīšu tikai divas atsevišķas ziņas precīzā formātā:} \\
\\ 
\texttt{Pirmā ziņa:} \\
\texttt{```} \\
\texttt{SAKI: Mani objekti: H(5,2) T(1,4) I(3,4) C(4,5) W(7,7)} \\
\texttt{```} \\
\\ 
\texttt{Otrā ziņa (ja pirmā tiks pieņemta):} \\
\texttt{```} \\
\texttt{SAKI: Piedāvātais mērķis: C(1,1) H(2,2) I(3,1) T(4,1) W(5,1)} \\
\texttt{```} \\
\\ 
\texttt{Vai vēlaties, lai es sūtu tikai pirmo ziņu? (Mēs esam tuvu maksimālajam sodu skaitam)} \\
            }
        }
    }
    & & \\ \\

    \theutterance \stepcounter{utterance}  
    & & & \multicolumn{2}{p{0.3\linewidth}}{
        \cellcolor[rgb]{0.9,0.9,0.9}{
            \makecell[{{p{\linewidth}}}]{
                \texttt{\tiny{[GM$|$GM]}}
                \texttt{Invalid format: head and tail are not empty} \\
\texttt{head: 'Atvainojos par iepriekšējām kļūdām. Tagad es sūtīšu tikai divas atsevišķas ziņas precīzā formātā:} \\
\\ 
\texttt{Pirmā ziņa:} \\
\\ 
\texttt{SAKI: Mani objekti: H(5,2) T(1,4) I(3,4) C(4,5) W(7,7)} \\
\\ 
\\ 
\texttt{Otrā ziņa (ja pirmā tiks pieņemta):} \\
\\ 
\texttt{'} \\
\texttt{tail: '} \\
\\ 
\\ 
\texttt{Vai vēlaties, lai es sūtu tikai pirmo ziņu? (Mēs esam tuvu maksimālajam sodu skaitam)'} \\
            }
        }
    }
    & & \\ \\

    \theutterance \stepcounter{utterance}  
    & & \multicolumn{4}{p{0.6\linewidth}}{
        \cellcolor[rgb]{0.9,0.9,0.9}{
            \makecell[{{p{\linewidth}}}]{
                \texttt{\tiny{[GM$\rangle$P2]}}
                \texttt{Sods: Jūsu ziņojumā nedrīkst būt nekas pirms vai pēc komandas!} \\
\\ 
\texttt{Pārliecinieties, ka jūsu atbildē ir tikai `SAKI: <ZIŅA>` vai `GĀJIENS: <OBJEKTS>, (<X>, <Y>)`, un nekas cits!} \\
\\ 
\texttt{Jūs kopā esat uzkrājuši 11 no 12 sodiem.} \\
\texttt{Lūdzu, mēģiniet vēlreiz!} \\
            }
        }
    }
    & \\ \\

    \theutterance \stepcounter{utterance}  
    & & & \multicolumn{4}{p{0.6\linewidth}}{
        \cellcolor[rgb]{0.9,0.9,0.9}{
            \makecell[{{p{\linewidth}}}]{
                \texttt{\tiny{[GM$\langle$P2]}}
                \texttt{Tagad esmu pilnībā sapratis formātu. Šeit ir precīzi atbilstoša komanda:} \\
\\ 
\texttt{```} \\
\texttt{SAKI:H(5,2) T(1,4) I(3,4) C(4,5) W(7,7)} \\
\texttt{```} \\
            }
        }
    }
     \\ \\

    \theutterance \stepcounter{utterance}  
    & & & \multicolumn{2}{p{0.3\linewidth}}{
        \cellcolor[rgb]{0.9,0.9,0.9}{
            \makecell[{{p{\linewidth}}}]{
                \texttt{\tiny{[GM$|$GM]}}
                \texttt{Tagad esmu pilnībā sapratis formātu. Šeit ir precīzi atbilstoša komanda:} \\
\\ 
\texttt{```} \\
\texttt{SAKI:H(5,2) T(1,4) I(3,4) C(4,5) W(7,7)} \\
\texttt{```} \\
            }
        }
    }
    & & \\ \\

    \theutterance \stepcounter{utterance}  
    & & & \multicolumn{2}{p{0.3\linewidth}}{
        \cellcolor[rgb]{0.9,0.9,0.9}{
            \makecell[{{p{\linewidth}}}]{
                \texttt{\tiny{[GM$|$GM]}}
                \texttt{Invalid response format} \\
            }
        }
    }
    & & \\ \\

    \theutterance \stepcounter{utterance}  
    & & \multicolumn{4}{p{0.6\linewidth}}{
        \cellcolor[rgb]{0.9,0.9,0.9}{
            \makecell[{{p{\linewidth}}}]{
                \texttt{\tiny{[GM$\rangle$P2]}}
                \texttt{Sods: Jūsu ziņojums nav gaidītajā formātā!} \\
\\ 
\texttt{Pārliecinieties, ka jūsu atbildē ir tikai `SAKI: <ZIŅA>` vai `GĀJIENS: <OBJEKTS>, (<X>, <Y>)`, un nekas cits!} \\
\\ 
\texttt{Jūs kopā esat uzkrājuši 12 no 12 sodiem.} \\
\texttt{Lūdzu, mēģiniet vēlreiz!} \\
            }
        }
    }
    & \\ \\

    \theutterance \stepcounter{utterance}  
    & & & \multicolumn{4}{p{0.6\linewidth}}{
        \cellcolor[rgb]{0.9,0.9,0.9}{
            \makecell[{{p{\linewidth}}}]{
                \texttt{\tiny{[GM$\langle$P2]}}
                \texttt{Es pilnībā saprotu {-} mēs esam sasnieguši maksimālo sodu skaitu. Pēdējais mēģinājums ar absolūti precīzu formātu:} \\
\\ 
\texttt{```} \\
\texttt{SAKI:H(5,2),T(1,4),I(3,4),C(4,5),W(7,7)} \\
\texttt{```} \\
\\ 
\texttt{Šoreiz:} \\
\texttt{1. Tikai `SAKI:` komanda} \\
\texttt{2. Tikai objekti ar koordinātām} \\
\texttt{3. Komatu atdalīti objekti} \\
\texttt{4. Bez atstarpēm pirms vai pēc komandas} \\
\texttt{5. Bez papildu vārdu vai simbolu} \\
\\ 
\texttt{Ja šis arī nebūs pareizs, esmu gatavs sākt pārvietošanas gājienus, lai izvairītos no spēles zaudēšanas.} \\
            }
        }
    }
     \\ \\

    \theutterance \stepcounter{utterance}  
    & & & \multicolumn{2}{p{0.3\linewidth}}{
        \cellcolor[rgb]{0.9,0.9,0.9}{
            \makecell[{{p{\linewidth}}}]{
                \texttt{\tiny{[GM$|$GM]}}
                \texttt{Es pilnībā saprotu {-} mēs esam sasnieguši maksimālo sodu skaitu. Pēdējais mēģinājums ar absolūti precīzu formātu:} \\
\\ 
\texttt{```} \\
\texttt{SAKI:H(5,2),T(1,4),I(3,4),C(4,5),W(7,7)} \\
\texttt{```} \\
\\ 
\texttt{Šoreiz:} \\
\texttt{1. Tikai `SAKI:` komanda} \\
\texttt{2. Tikai objekti ar koordinātām} \\
\texttt{3. Komatu atdalīti objekti} \\
\texttt{4. Bez atstarpēm pirms vai pēc komandas} \\
\texttt{5. Bez papildu vārdu vai simbolu} \\
\\ 
\texttt{Ja šis arī nebūs pareizs, esmu gatavs sākt pārvietošanas gājienus, lai izvairītos no spēles zaudēšanas.} \\
            }
        }
    }
    & & \\ \\

    \theutterance \stepcounter{utterance}  
    & & & \multicolumn{2}{p{0.3\linewidth}}{
        \cellcolor[rgb]{0.9,0.9,0.9}{
            \makecell[{{p{\linewidth}}}]{
                \texttt{\tiny{[GM$|$GM]}}
                \texttt{Invalid format: head and tail are not empty} \\
\texttt{head: 'Es pilnībā saprotu {-} mēs esam sasnieguši maksimālo sodu skaitu. Pēdējais mēģinājums ar absolūti precīzu formātu:} \\
\\ 
\\ 
\texttt{SAKI:H(5,2),T(1,4),I(3,4),C(4,5),W(7,7)} \\
\\ 
\\ 
\texttt{Šoreiz:} \\
\texttt{1. Tikai '} \\
\texttt{tail: '} \\
\texttt{2. Tikai objekti ar koordinātām} \\
\texttt{3. Komatu atdalīti objekti} \\
\texttt{4. Bez atstarpēm pirms vai pēc komandas} \\
\texttt{5. Bez papildu vārdu vai simbolu} \\
\\ 
\texttt{Ja šis arī nebūs pareizs, esmu gatavs sākt pārvietošanas gājienus, lai izvairītos no spēles zaudēšanas.'} \\
            }
        }
    }
    & & \\ \\

    \theutterance \stepcounter{utterance}  
    & & & \multicolumn{2}{p{0.3\linewidth}}{
        \cellcolor[rgb]{0.9,0.9,0.9}{
            \makecell[{{p{\linewidth}}}]{
                \texttt{\tiny{[GM$|$GM]}}
                \texttt{Maximum number of penalties exceeded} \\
            }
        }
    }
    & & \\ \\

    \theutterance \stepcounter{utterance}  
    & & & \multicolumn{2}{p{0.3\linewidth}}{
        \cellcolor[rgb]{0.9,0.9,0.9}{
            \makecell[{{p{\linewidth}}}]{
                \texttt{\tiny{[GM$|$GM]}}
                \texttt{Initial states:} \\
\texttt{```} \\
\texttt{Player 1:} \\
\texttt{ 1234567} \\
\texttt{╔═╤═════╗} \\
\texttt{║H│◌◌◌TI║ 1} \\
\texttt{╟─┤◌◌◌◌◌║ 2} \\
\texttt{║◌└─┐◌◌◌║ 3} \\
\texttt{║◌◌◌├─┐◌║ 4} \\
\texttt{║C◌◌│◌│◌║ 5} \\
\texttt{╟─┬─┴─┼─╢ 6} \\
\texttt{║◌│W◌◌│◌║ 7} \\
\texttt{╚═╧═══╧═╝} \\
\\ 
\\ 
\texttt{Player 2:} \\
\texttt{ 1234567} \\
\texttt{╔═╤═════╗} \\
\texttt{║◌│◌◌◌◌◌║ 1} \\
\texttt{╟─┤◌◌◌H◌║ 2} \\
\texttt{║◌└─┐◌◌◌║ 3} \\
\texttt{║T◌I├─┐◌║ 4} \\
\texttt{║◌◌◌│C│◌║ 5} \\
\texttt{╟─┬─┴─┼─╢ 6} \\
\texttt{║◌│◌◌◌│W║ 7} \\
\texttt{╚═╧═══╧═╝} \\
\\ 
\texttt{```} \\
            }
        }
    }
    & & \\ \\

    \theutterance \stepcounter{utterance}  
    & & & \multicolumn{2}{p{0.3\linewidth}}{
        \cellcolor[rgb]{0.9,0.9,0.9}{
            \makecell[{{p{\linewidth}}}]{
                \texttt{\tiny{[GM$|$GM]}}
                \texttt{End states:} \\
\texttt{```} \\
\texttt{Player 1:} \\
\texttt{ 1234567} \\
\texttt{╔═╤═════╗} \\
\texttt{║H│◌◌◌TI║ 1} \\
\texttt{╟─┤◌◌◌◌◌║ 2} \\
\texttt{║◌└─┐◌◌◌║ 3} \\
\texttt{║◌◌◌├─┐◌║ 4} \\
\texttt{║C◌◌│◌│◌║ 5} \\
\texttt{╟─┬─┴─┼─╢ 6} \\
\texttt{║◌│W◌◌│◌║ 7} \\
\texttt{╚═╧═══╧═╝} \\
\\ 
\\ 
\texttt{Player 2:} \\
\texttt{ 1234567} \\
\texttt{╔═╤═════╗} \\
\texttt{║◌│◌◌◌◌◌║ 1} \\
\texttt{╟─┤◌◌◌H◌║ 2} \\
\texttt{║◌└─┐◌◌◌║ 3} \\
\texttt{║T◌I├─┐◌║ 4} \\
\texttt{║◌◌◌│C│◌║ 5} \\
\texttt{╟─┬─┴─┼─╢ 6} \\
\texttt{║◌│◌◌◌│W║ 7} \\
\texttt{╚═╧═══╧═╝} \\
\\ 
\texttt{```} \\
            }
        }
    }
    & & \\ \\

    \theutterance \stepcounter{utterance}  
    & & & \multicolumn{2}{p{0.3\linewidth}}{
        \cellcolor[rgb]{0.9,0.9,0.9}{
            \makecell[{{p{\linewidth}}}]{
                \texttt{\tiny{[GM$|$GM]}}
                \texttt{* success: False} \\
\texttt{* lose: False} \\
\texttt{* aborted: True} \\
\texttt{{-}{-}{-}{-}{-}{-}{-}} \\
\texttt{* Shifts: 0.00} \\
\texttt{* Max Shifts: 8.00} \\
\texttt{* Min Shifts: 4.00} \\
\texttt{* End Distance Sum: 23.93} \\
\texttt{* Init Distance Sum: 23.93} \\
\texttt{* Expected Distance Sum: 20.95} \\
\texttt{* Penalties: 13.00} \\
\texttt{* Max Penalties: 12.00} \\
\texttt{* Rounds: 0.00} \\
\texttt{* Max Rounds: 20.00} \\
\texttt{* Object Count: 5.00} \\
            }
        }
    }
    & & \\ \\

\end{supertabular}
}

\end{document}
